% !TEX root = ../main.tex
\chapter{Wprowadzenie do dziedziny}

Niniejszy rozdział wprowadza w dziedzinę, w której osadzona jest praca. Najpierw przedstawia generatywną sztuczną inteligencję: jej definicję, zasadę działania, zastosowania oraz ograniczenia. Następnie omawia duże modele językowe, czyli klasę modeli generatywnych, na której opiera się system MLLMS. Wprowadzone pojęcia są podstawą dla kolejnych rozdziałów, które opisują technologie, projekt systemu i wyniki badań.

\section{Generatywna sztuczna inteligencja}

\subsection{Definicja i miejsce generatywnej sztucznej inteligencji}

Sztuczna inteligencja (SI) obejmuje metody, które pozwalają komputerom wykonywać zadania wymagające zwykle ludzkiej inteligencji. Jej istotną częścią jest uczenie maszynowe (ang. Machine Learning, ML), w którym zależności nie są programowane wprost, lecz wyznaczane przez algorytm na podstawie danych. Szczególnym przypadkiem uczenia maszynowego są sztuczne sieci neuronowe (ang. Artificial Neural Networks), czyli modele zbudowane z warstw prostych jednostek obliczeniowych, których parametry dobierane są w procesie uczenia. Sieci składające się z wielu warstw określa się mianem uczenia głębokiego (ang. Deep Learning, DL). Umożliwiło ono modelowanie bardzo złożonych zależności w danych, takich jak tekst, obraz czy dźwięk.

Na uczeniu głębokim opiera się współczesna generatywna sztuczna inteligencja (ang. Generative Artificial Intelligence, GenAI). Według \cite{feuerriegel_generative_ai} są to techniki obliczeniowe zdolne do generowania pozornie nowych, sensownych treści, takich jak tekst, obrazy czy dźwięk, na podstawie danych treningowych. Modele generatywne nie ograniczają się więc do klasyfikowania danych wejściowych lub przewidywania pojedynczej wartości, lecz tworzą nowe dane, przypominające te, na których były uczone. Wśród nich szczególne miejsce zajmują duże modele językowe, opisane w sekcji \ref{sec:llm}.

\subsection{Zasada działania i zastosowania}

Modele generatywne uczą się rozkładu prawdopodobieństwa danych treningowych $p(x)$, a w wersji warunkowej rozkładu $p(x \mid c)$, w którym $c$ oznacza kontekst, na przykład polecenie użytkownika (ang. prompt). Generowanie polega na losowaniu (ang. sampling) kolejnych elementów z tego rozkładu. W modelach językowych elementami są tokeny, czyli fragmenty tekstu, a odpowiedź powstaje token po tokenie: każdy kolejny wybierany jest na podstawie zapytania i dotychczas wygenerowanego tekstu. Stopień losowości reguluje parametr zwany temperaturą (ang. temperature). Dla tego samego zapytania model może więc zwrócić różne odpowiedzi, a ich powtarzalność zapewnia dopiero ustalone ziarno generatora liczb losowych (ang. seed).

Modele dyskryminatywne uczą się natomiast przypisywać dane wejściowe do kategorii, czyli modelują $p(y \mid x)$, gdzie $y$ jest etykietą \cite{ng_jordan_discriminative_generative}. Typowym przykładem jest filtr spamu rozstrzygający, czy wiadomość jest niechcianą pocztą. Pytanie wielokrotnego wyboru wygląda jak zadanie klasyfikacji, jednak model generatywny nie wybiera klasy, lecz generuje tekst, w którym znajduje się litera odpowiedzi. Jego odpowiedź jest więc wynikiem losowania, a nie deterministycznej decyzji klasyfikatora.

Treści tworzone przez modele generatywne mogą mieć różną postać. Według \cite{feuerriegel_generative_ai} są to między innymi tekst, obrazy i wideo, mowa i muzyka oraz kod źródłowy. Niniejsza praca dotyczy wyłącznie tekstu.

Generatywna sztuczna inteligencja jest wykorzystywana w edukacji, nauce, medycynie, biznesie i codziennej pracy z informacją, między innymi do odpowiadania na pytania, streszczania i tłumaczenia tekstów, wyjaśniania zagadnień oraz pomocy w programowaniu. Szczególnie istotne dla tej pracy są zadania wymagające rozumowania i wiedzy specjalistycznej, takie jak rozwiązywanie zadań matematycznych i odpowiadanie na pytania medyczne, na których oparto badania opisane w dalszej części pracy. W tych dziedzinach wiarygodność odpowiedzi ma duże znaczenie, co prowadzi do ograniczeń omówionych w następnej podsekcji.

\subsection{Ograniczenia generatywnej sztucznej inteligencji}

Mimo szerokich możliwości generatywna sztuczna inteligencja ma ograniczenia, które wynikają bezpośrednio ze sposobu jej działania.

Najbardziej znanym z nich są halucynacje (ang. hallucinations), czyli generowanie treści brzmiących wiarygodnie, lecz nieprawdziwych. Wyróżnia się wśród nich halucynacje dotyczące faktów, gdy treść jest sprzeczna z weryfikowalną wiedzą o świecie, oraz halucynacje dotyczące wierności, gdy odpowiedź odbiega od polecenia lub kontekstu \cite{llm_hallucination_survey}. Ponieważ model losuje tokeny zgodnie z wyuczonym rozkładem i nie sprawdza prawdziwości treści, odpowiedź błędna wygląda tak samo jak poprawna, a przekonujący ton utrudnia wykrycie błędu. Ma to duże znaczenie w dziedzinach takich jak medycyna, w których błędna informacja może prowadzić do szkody.

Z tym wiąże się brak samoweryfikacji. Model nie posiada wbudowanego mechanizmu sprawdzania własnej odpowiedzi, a poproszenie go o ponowną ocenę własnego rozumowania nie musi prowadzić do poprawy. Pojedynczy model jest więc jednocześnie autorem i jedynym recenzentem odpowiedzi.

Kolejnym ograniczeniem jest wrażliwość na sformułowanie polecenia. Słaby prompt, czyli ogólnikowy lub pozbawiony kontekstu, zwykle daje odpowiedź niepełną lub błędną, natomiast dobrze zbudowany, precyzujący zadanie i oczekiwany format, poprawia wynik. Nawet drobne zmiany formy polecenia, niezmieniające jego sensu, mogą zmienić odpowiedź.

Ograniczenia te, w szczególności brak samoweryfikacji, stanowią punkt wyjścia dla koncepcji współpracy wielu modeli, opisanej w kolejnych rozdziałach.

\section{Duże modele językowe}

\subsection{Definicja i architektura}

Duże modele językowe to nic innego jak sieci neuronowe o bardzo dużej liczbie parametrów, trenowane na ogromnych zbiorach tekstów. Opierają się one na architekturze Transformer \cite{attention_is_all_you_need}, która umożliwia efektywne modelowanie zależności w długich sekwencjach. Dzięki temu mogą one rozpoznawać i interpretować złożone teksty. Przed ich pojawieniem się, do tego typu zadań stosowano rekurencyjne sieci neuronowe (RNN \cite{fudamentals_of_rnn_and_lstm_network}), które zaczęły mieć trudności wraz z zwiększeniem modeli i długości sekwencji. Przez potrzebę wykonywania sekwencyjnych obliczeń były one mało efektywne.

Pojawienie się obecnych modeli uznaje się za wyraźny przełom w przetwarzaniu języka naturalnego (NLP) i w całej dziedzinie sztucznej inteligencji (SI)\cite{ma_collaboration}. Odpowiedzialna za to jest wcześniej wspomniana architektura Transformer. Opublikowana w 2017 roku praca "Attention Is All You Need" \cite{attention_is_all_you_need} wprowadziła mechanizm samo-uwagi (ang. self-attention), który pozwala modelom dynamicznie ważyć znaczenie poszczególnych słów w kontekście całej sekwencji. Poza tym architektura ta umożliwia równoległe przetwarzanie danych, co znacząco zwiększa efektywność treningu i wnioskowania w porównaniu do wcześniejszych podejść opartych na RNN. Od tego momentu zaczęto tworzyć coraz to większe modele. Liczba parametrów współczesnych modeli wynosi miliardy, a w największych przypadkach sięga bilionów ($10^{12}$). Przełomem okazały się rodziny modeli takie jak GPT, LLaMA i Gemini, które pokazały, że wzrost rozmiaru, ilości danych i mocy obliczeniowej przekłada się na jakość działania \cite{ma_collaboration}.

\subsection{Uczenie dużych modeli językowych}

Trening dużego modelu językowego przebiega w dwóch głównych etapach: wstępnym uczeniu (ang. pre-training) i uczeniu końcowym (ang. post-training).

Podczas pre-trainingu model uczy się na dużych zbiorach tekstu bez etykiet, a jego zadanie jest jedno: przewidzieć kolejny token na podstawie poprzednich. Taki trening wystarcza, by model nauczył się budować zdania poprawne pod względem składni. Matematycznie model uczy się rozkładu prawdopodobieństwa warunkowego $P(w_t \mid w_1, \dots, w_{t-1})$, czyli prawdopodobieństwa kolejnego tokenu przy danym kontekście. Dopiero przy generowaniu trzeba z tego rozkładu wybrać token. Najprostsza strategia, zachłanne dekodowanie (ang. greedy decoding), w każdym kroku bierze token najbardziej prawdopodobny w danym kontekście:
\[
w_t = \arg\max_{w} P(w \mid w_1, w_2, \dots, w_{t-1})
\]
Jakość przewidywań mierzy funkcja straty, którą w trakcie uczenia minimalizujemy:
\[
L = -\sum_{t=1}^{T} \log P(w_t \mid w_1, w_2, \dots, w_{t-1})
\]

W efekcie sieć przyswaja zasady gramatyki i ogólną budowę języka. Powstaje model bazowy, który przewiduje najbardziej prawdopodobną kontynuację tekstu. To nie to samo co wykonanie polecenia użytkownika, bo cel uczenia (przewidzenie kolejnego tokenu) różni się od celu, jakim jest pomocne i bezpieczne wykonywanie poleceń \cite{training_language_models}.

Tę lukę zmniejsza post-training. W schemacie opisanym dla modelu InstructGPT składa się on z trzech kroków \cite{training_language_models}. Najpierw model dostraja się nadzorowanie (ang. supervised fine-tuning, SFT) na przykładach pożądanych odpowiedzi napisanych przez ludzi. Następnie ludzie oceniają kilka odpowiedzi modelu na to samo polecenie, a na tych porównaniach uczy się model nagrody (ang. reward model), który przewiduje, którą odpowiedź woli człowiek. W ostatnim kroku model dostraja się uczeniem ze wzmocnieniem tak, by maksymalizował tę nagrodę (ang. reinforcement learning from human feedback, RLHF; w InstructGPT użyto algorytmu PPO). Podobny trzyetapowy schemat opisują Feuerriegel i in. dla ChatGPT \cite{feuerriegel_generative_ai}.

Po takim uczeniu model potrafi wykonywać polecenia wydawane przez użytkownika. Po odpowiednim sformułowaniu prompta potrafi przetłumaczyć tekst, odpowiedzieć w narzuconym formacie albo przyjąć wskazaną rolę. Dwie ostatnie umiejętności wykorzystano w tej pracy.

Dane treningowe pochodzą przeważnie z internetu i zostały napisane przez ludzi. Stąd bierze się przypuszczenie, że modele przejmują również część wzorców, według których ludzie podejmują decyzje.

\subsection{Sterowanie, wnioskowanie i wyzwania wdrożeniowe}

\subsubsection{Inżynieria promptów}

Zachowanie dużego modelu językowego można w znacznej mierze ukierunkować bez zmiany jego wag. Zajmuje się tym inżynieria promptów (ang. prompt engineering): dobór instrukcji, kontekstu, formatu odpowiedzi i przykładów tak, by model dał wynik zgodny z oczekiwaniami. Kosztuje to znacznie mniej niż ponowne trenowanie lub dostrajanie modelu, bo wymaga tylko zmiany tekstu wejściowego. Dlatego prompting jest zwykle pierwszym narzędziem, po które sięga się przy dopasowywaniu modelu do konkretnego zadania. 

\subsubsection{Chain-of-Thought}

Chain-of-Thought (CoT) to technika, która zamiast prosić o samą odpowiedź, skłania się model do rozbicia złożonego problemu na logiczne kroki pośrednie i dopiero na ich podstawie sformułowania wyniku \cite{llm_mas_techniques_and_business}. Dzięki rozłożeniu zadania na mniejsze kroki zwiększa się czytelność i staranność rozumowania modelu. Model przestaje "zgadywać" końcowy wynik. Dzięki temu łatwiej radzi sobię z problemami wymagającymi kilku następujących po sobie operacji. Technika ta świetnie sprawdza się w arytmetyce, logice czy planowaniu. Jako efekt uboczny uzyskujemy też jawny zapis toku rozumowania, który można przejrzeć i ocenić.

Samo CoT ma jednak słabość: zwykle łączy się je z dekodowaniem zachłannym (ang. greedy decoding), w którym w każdym kroku wybierany jest najbardziej prawdopodobny token. Powstaje wtedy jedna ścieżka rozumowania i jeden błąd w jej środku przesądza o wyniku. Metoda Self-Consistency (CoT-SC) zastępuje ten mechanizm wielokrotnym próbkowaniem: model generuje wiele różnych ścieżek rozumowania, zamiast zatrzymać się na jednej, zachłannej\cite{more_agents_is_all_you_need}\cite{self_consistency}. Następnie wybierana jest odpowiedź najbardziej spójna, czyli taka, do której prowadzi najwięcej niezależnych ścieżek

Cenę za tę poprawę płaci się w obliczeniach. Każde dodatkowe próbkowanie to kolejne wywołanie modelu, więc koszt i czas rosną proporcjonalnie do liczby generowanych ścieżek. W zastosowaniach lokalnych, gdzie zasoby sprzętowe są ograniczone, ma to realne znaczenie przy wyborze liczby próbek.

\subsubsection{Halucynacje}

Mimo imponujących możliwości duże modele językowe mają ograniczenia. Najczęściej wskazuje się wśród nich halucynacje, czyli generowanie błędnych lub niemających oparcia w faktach treści \cite{ma_collaboration}. Problem nie polega wyłącznie na tym, że model się myli. Błędne informacje podaje on tym samym płynnym, zdecydowanym tonem co poprawne, więc z samej formy odpowiedzi trudno ocenić, czy można jej zaufać. Wynika to ze sposobu działania modelu: wybiera on tekst najbardziej prawdopodobny w danym kontekście, a nie sprawdza go z wiarygodnym źródłem wiedzy.

Ryzyko rośnie w systemach złożonych z kilku współpracujących modeli lub agentów. Halucynacja jednego z nich może wywołać efekt kaskadowy: kolejne elementy systemu przyjmują błędny wynik jako pewnik i budują na nim dalsze rozumowanie, przez co błąd się powiela i wzmacnia \cite{llm_hallucination_survey}. Im dłuższy łańcuch przetwarzania, tym trudniej wskazać, w którym miejscu pojawił się pierwotny fałsz.

\label{sec:llm}

\section{Mechanizmy współpracy}

Ograniczenia opisane w poprzednich sekcjach, przede wszystkim halucynacje i brak samoweryfikacji, skłoniły badaczy do łączenia kilku wywołań modeli językowych w jeden system. Zamiast polegać na pojedynczej odpowiedzi, taki system zbiera kilka odpowiedzi, a następnie je porównuje, ocenia albo uzgadnia. Uczestnikami mogą być kolejne wywołania tego samego modelu albo różne modele. Każde dodatkowe wywołanie zwiększa koszt obliczeniowy, więc poszczególne mechanizmy różnią się nie tylko skutecznością, lecz także ceną.

\subsection{Głosowanie}

Najprostszym mechanizmem jest głosowanie większościowe. Kilka modeli albo kilka wywołań tego samego modelu rozwiązuje zadanie niezależnie, a za wynik przyjmuje się odpowiedź wskazaną najczęściej. Jest to ta sama idea co w metodzie Self-Consistency, rozszerzona na osobne wywołania i różne modele. W badaniach tej metody dokładność rosła wraz z liczbą agentów biorących udział w głosowaniu \cite{more_agents_is_all_you_need}. Modele nie wymieniają tu jednak informacji. Każdy pracuje w izolacji, a mechanizm jedynie zlicza odpowiedzi, więc błąd popełniony przez większość zostaje utrwalony. Przy dwóch uczestnikach głosowanie traci sens, bo w razie rozbieżności nie ma większości.

\subsection{Weryfikacja odpowiedzi przez drugi model}

Drugi mechanizm rozdziela role. Jeden model generuje rozwiązanie, a drugi je sprawdza: ocenia poprawność, wskazuje błędy i ewentualnie proponuje poprawkę. W ten sposób autor odpowiedzi przestaje być jej jedynym recenzentem. Ma to znaczenie, bo w badaniach modele proszone o poprawienie własnego rozumowania bez zewnętrznej informacji zwrotnej zwykle nie poprawiały wyniku, a niekiedy go pogarszały \cite{llm_cannot_self_correct}. Słabością tego schematu jest asymetria. Weryfikator też może się mylić, a jego ocena nie podlega kontroli, więc błędna krytyka może zepsuć poprawną odpowiedź.

\subsection{Debata}

Debata usuwa tę asymetrię. W jej podstawowym schemacie kilka modeli najpierw rozwiązuje zadanie niezależnie, a następnie w kolejnych rundach otrzymuje odpowiedzi pozostałych uczestników i aktualizuje własną \cite{multiagent_debate}. Każdy model jest więc jednocześnie autorem i recenzentem: broni swojego rozwiązania albo zmienia je pod wpływem argumentów pozostałych. Debata kończy się po ustalonej liczbie rund lub wcześniej, gdy odpowiedzi staną się zgodne. W badaniach schemat ten poprawił wyniki w zadaniach arytmetycznych i matematycznych oraz ograniczył halucynacje w pytaniach o fakty.

Skuteczność debaty zależy od różnorodności uczestników. Model, który nabrał pewności co do swojej odpowiedzi, przestaje generować nowe argumenty, nawet gdy odpowiedź jest błędna. Zjawisko to nazwano degeneracją myśli (ang. Degeneration-of-Thought) \cite{mad_divergent_thinking}. Kilka instancji tego samego modelu dzieli te same luki w wiedzy, więc łatwo osiągają zgodność także w błędzie. Różne modele, trenowane na innych danych, częściej mylą się w różnych miejscach, a wtedy rozbieżność ich odpowiedzi sygnalizuje, że zadanie wymaga dalszej dyskusji. Debata nie chroni jednak przed efektem kaskadowym opisanym przy halucynacjach. Jeśli jeden model przekonująco przedstawi błędne rozumowanie, drugi może je przejąć, a zgodność odpowiedzi nie gwarantuje ich poprawności.
